% Terjemahan Bahasa Indonesia dari chapter2.tex baris 9--32.
% Sumber: Methods of Algebra, Volume 2, commit
% 9a5803ff77dd3257484cb177f851a73770a59dd3 (CC BY 4.0).
% stable-unit-id: o014.aljabr2.chapter2.overview

% segment-id: o014.li2.chapter2.overview.h001
\chapter{Kategori Abelian}\label{sec:Abel-cat}
\hypertarget{o014.aljabr2.chapter2.overview}{ }

% segment-id: o014.li2.chapter2.overview.p001
Kategori abelian adalah kategori aditif yang mempunyai kernel dan kokernel,
serta semua morfismenya ketat. Untuk suatu gelanggang $R$, kategori ini
memperumum kategori modul kiri $R\dcate{Mod}$ atau kategori modul kanan
$\cated{Mod}R$. Di dalamnya dapat dikaji beragam konsep, seperti objek
injektif, objek projektif, kompleks, homologi/kohomologi, dan barisan eksak.

% segment-id: o014.li2.chapter2.overview.p002
Kategori modul merupakan purwarupa kategori abelian, tetapi belum menggambarkan
keseluruhan hakikatnya. Ada dua alasan.
% segment-id: o014.li2.chapter2.overview.l001
\begin{enumerate}
% segment-id: o014.li2.chapter2.overview.i001
	\item Dalam kategori abelian pada umumnya, objek tidak mempunyai elemen dan
	morfisme bukan pemetaan. Karena itu, metode pengejaran diagram dari kategori
	modul tidak lagi dapat digunakan; barisan-barisan eksak paling mendasar
	harus diturunkan melalui operasi pada objek dan funktor. Contoh yang
	menonjol ialah lema ular dan lema lima yang dibahas dalam
	\S~\sourcecrossref{sec:diagram-lemmas}{2.3}; pembuktiannya bertumpu pada
	semacam pengganti pengejaran diagram
	(Lema~\sourcecrossref{prop:Abel-cat-exact-aux}{pada \S~2.3}). Meskipun
	demikian, argumennya tetap jauh lebih rumit daripada versi untuk kategori
	modul.

% segment-id: o014.li2.chapter2.overview.i002
	\item Bahkan jika kita berangkat dari kategori modul atau subkategori
	abeliannya, kategori abelian abstrak tetap muncul ketika kita mengkaji
	kategori funktor, lokalisasi, atau kuosien Serre. Dalam geometri dan
	topologi, kajian teori berkas secara alami mengarah pada konstruksi semacam
	ini.
\end{enumerate}

% segment-id: o014.li2.chapter2.overview.p003
Teorema pembenaman Freyd--Mitchell
\sourcecrossref{prop:FM}{dalam Lampiran~B.6} menyatakan bahwa setiap kategori
abelian kecil dapat dibenamkan secara penuh dan setia ke suatu
$\cated{Mod}R$, dengan funktor pembenaman yang mempertahankan barisan eksak.
Setelah hasil ini diterima, banyak pernyataan mengenai barisan eksak dapat
dengan mudah direduksi ke kasus $\cated{Mod}R$. Teknik reduksi tersebut
memungkinkan pengejaran diagram tetap digunakan, tetapi manfaat utamanya
mungkin berupa ketenangan psikologis: di satu pihak, reduksi itu tidak
menyentuh hakikat kategori abelian; di pihak lain, bagaimana pun teorema
pembenaman dibuktikan,
sejumlah fakta inti mengenai kategori abelian tidak dapat dielakkan dan harus
dibangun dari awal. Isi bab ini tidak bergantung pada teorema pembenaman dan,
secara logis, juga tidak bergantung pada teori modul. Namun, pembaca diharapkan
mempunyai pemahaman dasar tentang aljabar homologis pada kategori modul;
keterampilan mengejar diagram tetap berguna dan, dari segi urutan belajar,
mungkin masih diperlukan. Pembaca disarankan merujuk pada \cite{Li1} atau
bagian terkait dalam buku ajar lain. \S~\sourcecrossref{sec:FM}{B.6} pada
lampiran buku ini akan menurunkan teorema pembenaman sebagai suatu penerapan
konstruksi Ind.

% segment-id: o014.li2.chapter2.overview.p004
Kategori abelian yang kolengkap, mempunyai generator, serta memiliki
$\varinjlim$ kecil terfilter yang semuanya eksak disebut kategori Grothendieck.
Inilah pokok bahasan bagian terakhir,
\S~\sourcecrossref{sec:Grothendieck-cat}{2.10}; syarat-syarat tersebut
berkaitan langsung dengan asumsi mengenai ukuran himpunan. Kategori
Grothendieck kerap muncul dalam geometri dan mempunyai sifat-sifat yang lebih
baik, misalnya kelengkapan
(Korolari~\sourcecrossref{prop:Grothendieck-cat-complete}{pada \S~2.10}) dan
keberadaan resolusi injektif kanonik
(Teorema~\sourcecrossref{prop:Grothendieck-injective}{pada \S~2.10}).
Dibandingkan dengan kategori abelian umum, kategori Grothendieck lebih
menyerupai kategori modul; makna konkret pernyataan ini
akan dijelaskan oleh Teorema Gabriel--Popescu
\sourcecrossref{prop:GP}{dalam Lampiran~A.3}.

% segment-id: o014.li2.chapter2.overview.p005
Beberapa teorema isomorfisme dasar dalam teori modul tetap berlaku bagi
kategori abelian umum. Hal yang sama berlaku bagi konsep-konsep dasar seperti
subobjek, objek kuosien, dekomposisi jumlah langsung, objek tak terurai, objek
sederhana, dan deret komposisi; semuanya merupakan pokok penting dalam aljabar.
Konsep-konsep ini akan dibahas dalam
\S\S~\sourcecrossref{sec:direct-sum}{2.5}--%
\sourcecrossref{sec:semisimple}{2.7} pada bab ini. Sebagai contoh, Teorema
Krull--Remak--Schmidt tentang dekomposisi jumlah langsung mempunyai perumuman
abstrak (Teorema~\sourcecrossref{prop:Krull-Schmidt-gen}{pada \S~2.5}).
Bahasa teori kisi berguna untuk tujuan ini; lihat penjelasan dalam
\S~\sourcecrossref{sec:lattices}{2.4}. Perlu ditekankan bahwa berbagai
pernyataan yang melibatkan jumlah langsung tak hingga sering kali baru berlaku
dalam kategori Grothendieck.

% segment-id: o014.li2.chapter2.overview.p006
\S~\sourcecrossref{sec:Serre-subcat}{2.9} dalam bab ini dimulai dengan
definisi dan karakterisasi subkategori abelian, lalu memperkenalkan pembentukan
kuosien suatu kategori abelian terhadap jenis subkategori abelian tertentu,
yang disebut subkategori Serre; hasilnya disebut kuosien Serre.
Kuosien ini dicirikan oleh sifat universal dan dapat dikonstruksi melalui
lokalisasi yang diperkenalkan dalam \S\ref{sec:cat-localization}. Bagian
tersebut juga memperkenalkan grup $\mathrm{K}_0$ dari suatu kategori abelian
beserta sifat-sifat dasarnya, termasuk prinsip Euler--Poincaré
(Teorema~\sourcecrossref{prop:EP}{pada \S~2.9}) dan barisan eksak bagi
$\mathrm{K}_0$ dari kuosien Serre
(Proposisi~\sourcecrossref{prop:K0-exact-seq}{pada \S~2.9}). Teknik-teknik ini
lazim dalam penerapan. Perlu dicatat bahwa konstruksi grup $\mathrm{K}_0$
berlaku bagi kelas struktur yang lebih luas daripada kategori abelian, yaitu
\emph{kategori eksak}; teorinya dapat dipelajari dalam \cite{Bu10}. Grup
$\mathrm{K}_0$ sendiri hanyalah titik awal dari deret grup abelian
$\mathrm{K}_n$ ($n\in\Z_{\geq 0}$). Pemahaman yang memadai tentang grup K
tingkat tinggi memerlukan sudut pandang teori homotopi, terutama konsep
``spektrum'', dan berada di luar cakupan buku ini.
\index{kategori eksak (exact category)}

% segment-id: o014.li2.chapter2.overview.n001
\begin{petunjukbacaan}
	Sebagaimana disebutkan di atas, pembaca sebaiknya memahami sampai tingkat
	tertentu operasi konkret dalam kategori modul. Untuk teori dasar kompleks
	beserta kohomologinya, materi
	\S\S~\sourcecrossref{sec:Abel-cat-def}{2.1}--%
	\sourcecrossref{sec:inj-proj}{2.8} wajib dikuasai. Selain pembahasan
	subkategori abelian, isi \S~\sourcecrossref{sec:Serre-subcat}{2.9} relatif
	jarang digunakan dalam bagian-bagian berikutnya, tetapi tetap merupakan
	pengetahuan umum dan berkaitan dengan sejumlah definisi dalam pembahasan
	kategori turunan. \S~\sourcecrossref{sec:Grothendieck-cat}{2.10} tentang
	kategori Grothendieck merupakan materi tambahan yang memerlukan argumen
	relatif rumit dan pengetahuan dari lampiran; pembaca pemula dapat
	melewatinya terlebih dahulu.

	Kompleks yang dinyatakan sebagai data $(X^n,d^n)_n$ (dengan
	$d^n:X^n\to X^{n+1}$) juga disebut kompleks kokrantai. Jika alih-alih itu
	digunakan indeks bawah, yakni $(X_n,d_n)_n$ dengan
	$d_n:X_n\to X_{n-1}$, kompleks itu disebut kompleks rantai. Selain dalam
	Bab~\sourcecrossref{sec:simplicial}{8}, buku ini terutama menggunakan
	notasi kompleks kokrantai.
\end{petunjukbacaan}
